% =============================================================
%  American Terawatt — SMP Level 1, single-line diagram
%
%  Structure:
%    smp-level-1.tex             <- this file
%    theme -> ../theme           <- symlink: theme/atstyle.sty and theme/fonts/
%                                   resolve as in the main presentation
%    .latexmkrc                  <- adds ~/hids/template/glyphwright to TEXINPUTS
%    figures/figures.yaml        <- figure list; figures/smp_level_1.tex is
%    figures/src/smp_level_1.yaml   generated from the diagram YAML
%
%  Compile in this folder:
%    latexmk -xelatex smp-level-1.tex
%
%  The symbols come from the Glyphwright library in ~/hids/template/glyphwright
%  (\usepackage{glyphwright}); regenerate figure and library as described in
%  figures/figures.yaml.
% =============================================================
\documentclass[aspectratio=169,11pt]{beamer}
\usepackage{theme/atstyle}
\usepackage{glyphwright}

\title{HVDC Design}
\subtitle{AmWatt Team Presenation}
\author{American Terawatt}
\date{\today}

\newcommand{\slfig}[1]{\input{figures/#1}}

\tikzset{pboth/.style={{Latex[length=2mm]}-{Latex[length=2mm]}},
         pac2dc/.style={-{Latex[length=2mm]}},
         pdc2ac/.style={{Latex[length=2mm]}-}}

\begin{document}

\titleslide

% ---- Single-line diagram ------------------------------------
\begin{frame}{System Overview}
  ORC generators feed a data centre through a point-to-point VSC HVDC link.
  \vspace{2.2em}

  \centering
  
  \sbox0{\slfig{smp_level_1_gen_load}}
  
  \begin{tikzpicture}
    \node[inner sep=0] (fig) {\usebox0};
    \draw[-{Latex[length=2mm]}, violet, line width=1.2pt]
      ([xshift=-6mm]fig.center) -- node[above] {$P$} ([xshift=6mm]fig.center);
  \end{tikzpicture}
\end{frame}

\begin{frame}{System Overview}
  The complete complete system consts of different elements, supplied by different companies.
  AmWatt is reponsible for the integration of those elements.
  
  \vspace{2.2em}

  \centering
  % Overlays on the generated figure (kept out of figures/):
  % section boxes in the figure's own coordinates, x from/to and colour:
  % generation, busbar, converter station, d.c. line, converter station,
  % busbar, data centre; then the power-flow arrow on top
  \sbox0{\tikzset{every picture/.append style={execute at end picture={
      \foreach \xa/\xb/\c in {-51/-20.5/atLime, -20.5/-14.5/atGray,
                             -14.5/52/cyan, 52/95/violet, 95/162/cyan,
                             162/168/atGray, 168/193/orange}
        \draw[\c, fill=\c, fill opacity=0.15, rounded corners=1pt]
          (\xa+0.4,-24) rectangle (\xb-0.4,24);
    }}}\slfig{smp_level_1_gen_load}}%
  \begin{tikzpicture}
    \node[inner sep=0] (fig) {\usebox0};
    \draw[-{Latex[length=2mm]}, violet, line width=1.2pt]
      ([xshift=-6mm]fig.center) -- node[above] {$P$} ([xshift=6mm]fig.center);
  \end{tikzpicture}
\end{frame}


% ---- VSC converter -------------------------------------------
\begin{frame}{VSC converter}
  \noindent The VSC converts the a.c.\ voltage $U_{ac}$ into the pole-to-pole d.c.\ voltage $U_{dc}$.
  \vspace{1.5em}

  \centering
  \begin{tikzpicture}[x=1mm, y=1mm, scale=0.7, transform shape,
      volt/.style={gw/thin, gw/c-voltage, -{Latex[length=2mm]}}]
    \pic[gw/signals={none}, gw/no label] (U) at (0,0) {gw/vsc2-sig};
    % a.c. connection with voltage to earth
    \draw[gw/normal] (U-ac) -- ++(-12,0) coordinate (ac);
    \fill (ac) circle (0.5);
    \draw[volt] ($(ac)+(0,-0.75)$) -- ++(0,-16.25);
    \foreach \y/\w in {-18.25/2.5, -19.5/1.5, -20.75/0.5}
      \draw[gw/normal] ($(ac)+(-\w,\y)$) -- ++(2*\w,0);
    \node[gw/c-voltage, anchor=east] at ($(ac)+(-1.25,-9)$) {$U_{ac}$};
    % d.c. terminals with pole-to-pole voltage
    \fill (U-dc1) circle (0.5) (U-dc2) circle (0.5);
    \draw[volt] ($(U-dc1)+(0,-0.75)$) -- ($(U-dc2)+(0,0.75)$);
    \node[gw/c-voltage, anchor=west] at ($(U-dc1)!0.5!(U-dc2)+(1.25,0)$) {$U_{dc}$};
  \end{tikzpicture}
\end{frame}

% ---- VSC operating modes -------------------------------------
\begin{frame}{Converter operating modes}
  \noindent The same converter works as rectifier or inverter, set by the direction of $P$.
  \vspace{0.3em}

  \centering
  \begin{tikzpicture}[x=1mm, y=1mm, scale=0.5,
      pflow/.style={violet, line width=1.2pt},
      mode/.style={font=\small\bfseries}]
    \foreach \n/\x/\y/\name/\dir in {C/0/0/Converter/pboth,
                                     I/-55/-65/Inverter/pdc2ac,
                                     R/55/-65/Rectifier/pac2dc} {
      \pic[transform shape, gw/signals={none}, gw/no label] (\n) at (\x,\y) {gw/vsc2-sig};
      \node[mode, anchor=south] at (\x+5,\y+23) {\name};
      \draw[pflow, \dir] (\x-15,\y-27) -- node[below] {$P$} (\x+25,\y-27);
    }
    \node[font=\scriptsize, color=atGray, anchor=east] at (C-ac) {a.c.};
    \node[font=\scriptsize, color=atGray, anchor=west] at ($(C-dc1)!0.5!(C-dc2)$) {d.c.};
  \end{tikzpicture}
\end{frame}


% \begin{frame}{System Overview}
%   ORC generators feed a data centre through a point-to-point VSC HVDC link.
%   \vspace{2.2em}

%   \centering
  
%   \sbox0{\slfig{smp_level_1}}
  
%   \begin{tikzpicture}
%     \node[inner sep=0] (fig) {\usebox0};
%     \draw[-{Latex[length=2mm]}, violet, line width=1.2pt]
%       ([xshift=-6mm]fig.center) -- node[above] {$P$} ([xshift=6mm]fig.center);
%   \end{tikzpicture}
% \end{frame}

\end{document}

